\documentclass[11pt, letterpaper, twoside, openright]{article}
\usepackage{geometry}
\usepackage[utf8]{inputenc}
\usepackage{amsmath}
\usepackage{amsfonts}
\usepackage{amssymb}
\usepackage{graphicx}
\usepackage{hyperref}
\usepackage{fancyhdr}
\usepackage{geometry}
\usepackage{ragged2e}

\usepackage[T1]{fontenc}
\usepackage{accanthis}
%|-----------------|
%|  No mover nada  |
%|-----------------|

\geometry{
  left=2.5cm, % Márgen izquierdo
  right=2.5cm, % Márgen derecho
  top=3cm, % Márgen superior ajustado para incluir el encabezado
  bottom=4cm % Márgen inferior
}

% Configuración del encabezado
\pagestyle{fancy}
\fancyhf{}
\fancyhead[L]{\includegraphics[width=1.3cm]{logo.jpg}} 

\fancyhead[C]{\textbf{Unidad Profesional Interdisciplinaria de \\Ingeniería Campus Tlaxcala} \\ Dr. Jesús García Ramírez} % Cambia la materia
\fancyhead[R]{\includegraphics[width=3cm]{UPIIT.jpg}}
\fancyfoot[C]{\thepage}
\renewcommand{\headrulewidth}{1pt} % Línea bajo el encabezado
\renewcommand{\footrulewidth}{1pt}

% Ajuste del espacio entre el encabezado y el texto
\setlength{\headsep}{1.9cm} % Espacio entre el encabezado y el inicio del texto

% Aplicar encabezado a la primera página
\fancypagestyle{plain}{
  \fancyhf{}
  \fancyhead[L]{\includegraphics[width=1.3cm]{logo.jpg}}
  \fancyhead[C]{\textbf{Unidad Profesional Interdisciplinaria de \\Ingeniería Campus Tlaxcala} \\ Dr. Jesús García Ramírez \\ Materia} % Cambia este texto por el que desees
  \fancyhead[R]{\includegraphics[width=3cm]{UPIIT.jpg}}
  \fancyfoot[C]{\thepage}
  \renewcommand{\headrulewidth}{1pt}
  \renewcommand{\footrulewidth}{1pt}
}
\renewcommand{\refname}{Referencias}
%|-----------------|
%|  No mover nada  |
%|-----------------|


\title{Tarea \#1}
\author{Nombre del Estudiante(s)}
\date{}



\begin{document}





\vspace*{10pt}
\centering
\huge
Nombre de la materia

\LARGE
\textbf{Título de la Tarea}

\Large
Nombre del estudiante

Nombre del estudiante

\normalsize
\justifying

\begin{abstract}
Este es el resumen del documento. Aquí se debe proporcionar un breve resumen del contenido y los objetivos del documento.
%\keywords{Palabras clave, Separadas, Por comas}
\end{abstract}

%----------------------------------
% Para reportes finales, incluir la tabla de contenidos
%\tableofcontents
%\newpage
%----------------------------------

\section{Introducción}
Escribe aquí la introducción de tu documento. Este documento sigue un estilo parecido a la serie Lecture Notes in Computer Science (LNCS).

\section{Trabajo Relacionado}
Describe aquí el trabajo relacionado con el tema de tu documento. Puedes citar trabajos previos usando comandos como \cite{lamport1994latex}.

\section{Metodología}
Detalla la metodología que has seguido para llevar a cabo tu investigación o estudio.

\section{Resultados}
Presenta aquí los resultados obtenidos. Puedes incluir figuras, tablas y ecuaciones.

\subsection{Tabla}
Y aquí incluimos una tabla como ejemplo:

\begin{table}[htbp]
  \centering
  \caption{Ejemplo de tabla}
  \label{tab:ejemplo}
  \begin{tabular}{|c|c|}
    \hline
    \textbf{Columna 1} & \textbf{Columna 2} \\
    \hline
    Valor 1 & Valor A \\
    Valor 2 & Valor B \\
    \hline
  \end{tabular}
\end{table}

\subsection{Figura}
Aquí incluimos una figura como ejemplo:

\begin{figure}[htbp]
  \centering
  \includegraphics[width=5cm]{ejemplo.png} % Cambia 'ejemplo_figura.png' por el nombre de tu archivo de imagen
  \caption{Descripción de la figura.}
  \label{fig:ejemplo}
\end{figure}

\section{Conclusiones}
Escribe las conclusiones a las que has llegado con tu estudio o investigación. También puedes discutir posibles trabajos futuros~\cite{knuth1997art}.

\section*{Agradecimientos}
Incluye los agradecimientos, si es necesario.
\bibliographystyle{splncs04}
\bibliography{referencias} % El archivo .bib debe llamarse 'referencias.bib'

\end{document}
